\documentclass[../../main.tex]{subfiles}

\begin{document}

\section{度量空间上的 Carath\'eodory 外测度与 Hausdorff 测度}\label{section:ch20-hausdorff}

欧氏空间 $\mathbf{R}^{n}$ 上的 Lebesgue 外测度具有以下性质: 若 $A$ 和 $B$ 是 $\mathbf{R}^{n}$ 的子集且存在 $\delta>0$ 使得对所有的 $u\in A$ 与 $v\in B$, 满足 $\|u-v\|\geqslant\delta$, 则
\begin{align*}
\mu_{n}^{*}(A\cup B)=\mu_{n}^{*}(A)+\mu_{n}^{*}(B).
\end{align*}
我们在这简短的一节研究以下测度: 具有该性质的由度量空间上的外测度诱导的测度以及称为 Hausdorff 测度的一类这样的测度.

令 $X$ 为集合而 $\Gamma$ 是 $X$ 上的实值函数族. 了解在何种条件下外测度 $\mu^{*}$ 具有以下性质: $\Gamma$ 中的每个函数关于由 $\mu^{*}$ 通过 Carath\'eodory 构造法诱导的测度是可测的, 常常是有趣的. 我们给出一个充分性的准则. $X$ 的两个子集 $A$ 和 $B$ 称为被 $X$ 上的实值函数 $f$ 分离, 若存在实数 $a$ 和 $b$ 满足 $a<b$, 使得在 $A$ 上 $f\leqslant a$ 而在 $B$ 上 $f\geqslant b$.

\begin{proposition}\label{proposition:ch20-分离函数可测}
令 $\varphi$ 为集合 $X$ 上的实值函数而 $\mu^{*}:2^{X}\to[0,\infty]$ 是外测度, 它具有以下性质: 只要 $X$ 的两个子集 $A$ 和 $B$ 可被 $\varphi$ 分离, 就有
\begin{align*}
\mu^{*}(A\cup B)=\mu^{*}(A)+\mu^{*}(B).
\end{align*}
则 $\varphi$ 关于由 $\mu^{*}$ 诱导的测度是可测的.
\end{proposition}

\begin{proof}
令 $a$ 为实数. 我们必须证明集合
\begin{align*}
E=\{x\in X\mid\varphi(x)>a\}
\end{align*}
是 $\mu^{*}$ 可测的, 即对任何 $\varepsilon>0$ 与 $X$ 的任何有限外测度的子集 $A$,
\begin{align}\label{align:ch20-可测性不等式}
\mu^{*}(A)+\varepsilon\geqslant\mu^{*}(A\cap E)+\mu^{*}(A\cap E^{C}).
\end{align}
定义 $B=A\cap E$ 与 $C=A\cap E^{C}$. 对每个自然数 $n$, 定义
\begin{align*}
B_{n}=\{x\in B\mid\varphi(x)>a+1/n\}\text{ 与 }R_{n}=B_{n}\setminus B_{n-1}.
\end{align*}
我们有
\begin{align*}
B=B_{n}\cup\left[\bigcup_{k=n+1}^{\infty}R_{k}\right].
\end{align*}
现在在 $B_{n-2}$ 上我们有 $\varphi>a+1/(n-2)$, 而在 $R_{n}$ 上我们有 $\varphi\leqslant a+1/(n-1)$. 因此 $\varphi$ 分离 $R_{n}$ 与 $B_{n-2}$, 于是分离 $R_{2k}$ 与 $\bigcup_{j=1}^{k-1}R_{2j}$. 由于后者包含于 $B_{2k-2}$, 因此, 我们用归纳法证明对每个 $k$,
\begin{align*}
\mu^{*}\left[\bigcup_{j=1}^{k}R_{2j}\right]=\mu^{*}(R_{2k})+\mu^{*}\left[\bigcup_{j=1}^{k-1}R_{2j}\right]=\sum_{j=1}^{k}\mu^{*}(R_{2j}).
\end{align*}
由于 $\sum_{j=1}^{k}R_{2j}\subseteq B\subseteq A$, 我们有 $\sum_{j=1}^{k}\mu^{*}(R_{2j})\leqslant\mu^{*}(A)$, 从而级数 $\sum_{k=1}^{\infty}\mu^{*}(R_{2j})$ 收敛. 类似地, 级数 $\sum_{j=1}^{\infty}\mu^{*}(R_{2j+1})$ 收敛, 且因此级数 $\sum_{k=1}^{\infty}\mu^{*}(R_{k})$ 收敛. 选取 $n$ 充分大使得 $\sum_{k=n+1}^{\infty}\mu^{*}(R_{k})<\varepsilon$. 则根据 $\mu^{*}$ 的可数单调性,
\begin{align*}
\mu^{*}(B)\leqslant\mu^{*}(B_{n})+\sum_{k=n+1}^{\infty}\mu^{*}(R_{k})<\mu^{*}(B_{n})+\varepsilon.
\end{align*}
或
\begin{align*}
\mu^{*}(B_{n})>\mu^{*}(B)-\varepsilon.
\end{align*}
现在
\begin{align*}
\mu^{*}(A)\geqslant\mu^{*}(B_{n}\cup C)=\mu^{*}(B_{n})+\mu^{*}(C).
\end{align*}
因为 $\varphi$ 分离 $B_{n}$ 与 $C$. 因此,
\begin{align*}
\mu^{*}(A)\geqslant\mu^{*}(B)+\mu^{*}(C)-\varepsilon.
\end{align*}
我们证明了所要的不等式\eqref{align:ch20-可测性不等式}.
\end{proof}

令 $(X,\rho)$ 为度量空间. 回忆对 $X$ 的两个子集 $A$ 和 $B$, 我们定义 $A$ 和 $B$ 的距离 $\rho(A,B)$ 为
\begin{align*}
\rho(A,B)=\inf_{u\in A,v\in B}\rho(u,v).
\end{align*}
谈到与该度量空间相应的 Borel $\sigma$ 代数, 记为 $\mathcal{B}(X)$, 我们指的是包含由该度量诱导的拓扑的最小 $\sigma$ 代数.

\begin{definition}[Carath\'eodory 外测度]\label{definition:ch20-caratheodory外测度}
令 $(X,\rho)$ 为度量空间. 外测度 $\mu^{*}:2^{X}\to[0,\infty]$ 称为 Carath\'eodory 外测度, 若只要 $A$ 和 $B$ 是 $X$ 的两个子集使得 $\rho(A,B)>0$, 就有
\begin{align*}
\mu^{*}(A\cup B)=\mu^{*}(A)+\mu^{*}(B).
\end{align*}
\end{definition}

\begin{theorem}\label{theorem:ch20-borel可测}
令 $\mu^{*}$ 为度量空间 $(X,\rho)$ 上的 Carath\'eodory 外测度. 则 $X$ 的每个 Borel 子集关于 $\mu^{*}$ 是可测的.
\end{theorem}

\begin{proof}
Borel 集族是包含闭集的最小 $\sigma$ 代数且可测集是 $\sigma$ 代数. 因此仅须证明每个闭集是可测的. 然而, $X$ 的每个闭子集 $F$ 可表示为 $F=f^{-1}(0)$, 其中 $f$ 是 $X$ 上的连续函数, 定义为 $f(x)=\rho(F,\{x\})$. 因此仅须证明每个连续函数是可测的. 为此, 我们运用命题\cref{proposition:ch20-分离函数可测}. 事实上, 令 $A$ 和 $B$ 为 $X$ 的子集使得存在 $X$ 上的连续函数和实数 $a<b$ 满足在 $A$ 上 $f\leqslant a$ 而在 $B$ 上 $f\geqslant b$. 根据 $f$ 的连续性, $\rho(A,B)>0$. 因此, 根据假设, $\mu^{*}(A\cup B)=\mu^{*}(A)+\mu^{*}(B)$. 根据命题\cref{proposition:ch20-分离函数可测}, 每个连续函数是可测的. 证明完毕.
\end{proof}

我们现在将注意力转向度量空间 $(X,\rho)$ 上特定的外测度族. 首先回忆我们定义 $X$ 的子集 $A$ 的直径 $\operatorname{diam}(A)$ 为
\begin{align*}
\operatorname{diam}(A)=\sup_{u,v\in A}\rho(u,v).
\end{align*}
固定 $\alpha>0$. 对每个正实数 $\alpha$, 我们定义 Borel 代数 $\mathcal{B}(X)$ 上称为 $X$ 上的 $\alpha$ 维 Hausdorff 测度 $H_{\alpha}$. 这些测度对欧氏空间 $\mathbf{R}^{n}$ 特别重要, 在这种情形下它们给出了 $n$ 维 Lebesgue 测度为零的集合的尺寸的分级.

固定 $\alpha>0$. 取 $\varepsilon>0$ 且对 $X$ 的子集 $E$, 定义
\begin{align*}
H_{\alpha}^{(\varepsilon)}(E)=\inf\sum_{k=1}^{\infty}[\operatorname{diam}(A_{k})]^{\alpha}
\end{align*}
其中 $\{A_{k}\}_{k=1}^{\infty}$ 是覆盖 $E$ 的 $X$ 的子集的可数族, 且每个 $A_{k}$ 有小于 $\varepsilon$ 的直径. 观察到 $H_{\alpha}^{(\varepsilon)}$ 随着 $\varepsilon$ 减少而增加. 定义
\begin{align*}
H_{\alpha}^{*}(E)=\sup_{\varepsilon>0}H_{\alpha}^{(\varepsilon)}(E)=\lim_{\varepsilon\to 0}H_{\alpha}^{(\varepsilon)}(E).
\end{align*}

\begin{proposition}\label{proposition:ch20-hausdorff外测度}
令 $(X,\rho)$ 为度量空间而 $\alpha$ 是正实数. 则集函数 $H_{\alpha}^{*}:2^{X}\to[0,\infty]$ 是 Carath\'eodory 外测度.
\end{proposition}

\begin{proof}
容易证明 $H_{\alpha}^{*}$ 是 $2^{X}$ 上的可数单调集函数且 $H_{\alpha}^{*}(\varnothing)=0$. 因此 $H_{\alpha}^{*}$ 是 $2^{X}$ 上的外测度. 我们宣称它是 Carath\'eodory 外测度. 事实上, 令 $E$ 和 $F$ 为 $X$ 的两个子集, 满足 $\rho(E,F)>\delta$. 则只要 $\varepsilon<\delta$,
\begin{align*}
H_{\alpha}^{(\varepsilon)}(E\cup F)\geqslant H_{\alpha}^{(\varepsilon)}(E)+H_{\alpha}^{(\varepsilon)}(F).
\end{align*}
若 $\{A_{k}\}$ 是覆盖 $E\cup F$ 的可数集族, 每个的直径至多为 $\varepsilon$, 没有 $A_{k}$ 与 $E$ 和 $F$ 都有非空交. 取当 $\varepsilon\to 0$ 时的极限我们有
\begin{align*}
H_{\alpha}^{*}(E\cup F)\geqslant H_{\alpha}^{*}(E)+H_{\alpha}^{*}(F).
\end{align*}
\end{proof}

我们从定理\cref{theorem:ch20-borel可测}推出 $H_{\alpha}^{*}$ 诱导了包含 $X$ 的 Borel 子集的 $\sigma$ 代数上的一个测度. 我们记该测度在 $\mathcal{B}(X)$ 上的限制为 $H_{\alpha}$, 且称它为度量空间 $X$ 上的 Hausdorff $\alpha$ 维测度.

\begin{proposition}\label{proposition:ch20-hausdorff零测}
令 $(X,\rho)$ 为度量空间, $A$ 是 $X$ 的 Borel 子集, $\alpha$, $\beta$ 是正实数, 其中 $\alpha<\beta$. 若 $H_{\alpha}(A)<\infty$, 则 $H_{\beta}(A)=0$.
\end{proposition}

\begin{proof}
令 $\varepsilon>0$. 选取 $\{A_{k}\}_{k=1}^{\infty}$ 为 $A$ 的直径小于 $\varepsilon$ 的集合的覆盖, 使得
\begin{align*}
\sum_{k=1}^{\infty}[\operatorname{diam}(A_{k})]^{\alpha}\leqslant H_{\alpha}(A)+1.
\end{align*}
则
\begin{align*}
H_{\beta}^{(\varepsilon)}(A)\leqslant\sum_{k=1}^{\infty}[\operatorname{diam}(A_{k})]^{\beta}\leqslant\varepsilon^{\beta-\alpha}\cdot\sum_{k=1}^{\infty}[\operatorname{diam}(A_{k})]^{\alpha}\leqslant\varepsilon^{\beta-\alpha}\cdot[H_{\alpha}(A)+1].
\end{align*}
取 $\varepsilon\to 0$ 的极限得出 $H_{\beta}(A)=0$.
\end{proof}

对 $\mathbf{R}^{n}$ 的子集 $E$, 我们定义 $E$ 的 Hausdorff 维数 $\dim_{H}(E)$ 为
\begin{align*}
\dim_{H}(E)=\inf\{\beta\geqslant 0\mid H_{\beta}(E)=0\}.
\end{align*}

Hausdorff 测度对欧氏空间 $\mathbf{R}^{n}$ 特别有意义. 在 $n=1$ 的情形, $H_{1}$ 等于 Lebesgue 测度. 为看到这一点, 对 $I\subseteq\mathbf{R}$ 为区间. 给定 $\varepsilon>0$, 区间 $I$ 可表示为长度小于 $\varepsilon$ 的子区间的不交并, 且每个子区间的直径是它的长度. 因此 $H_{1}$ 与 Lebesgue 测度在实数区间的半环上一致. 于是, 根据这些测度的构造, 这些测度在 Borel 集上也一致. 因此 $H_{1}^{(\varepsilon)}$ 是 $E$ 的 Lebesgue 外测度. 对 $n>1$, $H_{n}$ 不等于 Lebesgue 测度(见习题 48), 但可证明它是常数乘以 $n$ 维 Lebesgue 测度(见习题 55). 从上面命题得出, 若 $A$ 是具有正 Lebesgue 测度的 $\mathbf{R}^{n}$ 的子集, 则 $\dim_{H}(A)=n$. 欧氏空间的子集的 Hausdorff 维数有许多具体的计算. 例如, 可以证明 Cantor 集的 Hausdorff 维数是 $\log 2/\log 3$. 关于 Hausdorff 测度的进一步结果, 包括 Hausdorff 维数的具体计算, 可以在 Yakov Pesin 的书《Dimension Theory and Dynamical System》[Pes98] 中找到.

\subsection*{习题}

\begin{exercise}\label{exercise:ch20-46}
证明在 Hausdorff 测度的定义中人们可取用开集或闭集的覆盖.
\end{exercise}

\begin{exercise}\label{exercise:ch20-47}
证明外 Hausdorff 测度 $H_{\alpha}^{*}$ 这一集函数是可数单调的.
\end{exercise}

\begin{exercise}\label{exercise:ch20-48}
在平面 $\mathbf{R}^{2}$ 中证明有界集可被有相同直径的球包含. 以此证明对 $\mathbf{R}^{2}$ 的有界子集 $A$, $H_{2}(A)\geqslant 4/\pi\cdot\mu_{2}(A)$, 其中 $\mu_{2}$ 是 $\mathbf{R}^{2}$ 上的 Lebesgue 测度.
\end{exercise}

\begin{exercise}\label{exercise:ch20-49}
令 $(X,\rho)$ 为度量空间而 $\alpha>0$. 对 $E\subseteq X$, 定义
\begin{align*}
H_{\alpha}'(E)=\inf\sum_{k=1}^{\infty}[\operatorname{diam}(A_{k})]^{\alpha}
\end{align*}
其中 $\{A_{k}\}_{k=1}^{\infty}$ 是覆盖 $E$ 的 $X$ 的子集的可数族: 在覆盖中对集合的直径没有限制. 比较集函数 $H_{\alpha}'$ 与 $H_{\alpha}$.
\end{exercise}

\begin{exercise}\label{exercise:ch20-50}
证明欧氏空间 $\mathbf{R}^{n}$ 上的每个 Hausdorff 测度 $H_{\alpha}$ 关于刚体运动不变.
\end{exercise}

\begin{exercise}\label{exercise:ch20-51}
直接证明: 若 $I$ 是 $\mathbf{R}^{n}$ 中的非平凡区间, 则 $H_{n}(I)>0$.
\end{exercise}

\begin{exercise}\label{exercise:ch20-52}
证明在任何度量空间, $H_{0}$ 是计数测度.
\end{exercise}

\begin{exercise}\label{exercise:ch20-53}
令 $[a,b]$ 为实数的闭有界区间且考虑定义为 $R=\{(x,y)\in\mathbf{R}^{2}\mid a\leqslant x\leqslant b,\ y=0\}$ 的平面的子集. 证明 $H_{2}(R)=0$. 接着证明 $H_{1}(R)=b-a$. 得出 $R$ 的 Hausdorff 维数是 $1$.
\end{exercise}

\begin{exercise}\label{exercise:ch20-54}
令 $f:[a,b]\to\mathbf{R}$ 为闭有界区间 $[a,b]$ 上的连续有界函数, 它在开区间 $(a,b)$ 上有连续有界导数. 考虑 $f$ 的图像 $G$ 为平面的子集. 证明 $H_{1}(G)=\int_{a}^{b}\sqrt{1+|f'(x)|^{2}}dx$.
\end{exercise}

\begin{exercise}\label{exercise:ch20-55}
令 $J$ 为 $\mathbf{R}^{n}$ 中的区间, 每条边的长度为 $1$. 定义 $\gamma_{n}=H_{n}(J)$. 证明: 若 $I$ 是 $\mathbf{R}^{n}$ 中的任何有界区间, 则 $H_{n}(I)=\gamma_{n}\cdot\mu_{n}(I)$. 用 Carath\'eodory-Hahn 定理的唯一性断言由此推出在 $\mathbf{R}^{n}$ 的 Borel 子集上 $H_{n}=\gamma_{n}\cdot\mu_{n}$.
\end{exercise}

\end{document}